\documentclass[11pt]{article}
\usepackage[T1]{fontenc}
\usepackage{lmodern,amsmath,amssymb,amsthm,booktabs}
\usepackage[margin=1in]{geometry}
\usepackage{microtype,xurl,array,longtable}
\usepackage[colorlinks=true,linkcolor=blue,urlcolor=blue,citecolor=blue]{hyperref}
\hypersetup{pdftitle={Could the class-number-one theorem have been proved in 1891?},
 pdfauthor={Research note prepared with Codex},
 pdfsubject={A reconstruction from Weber's modular theory with a historical source audit}}
\newtheorem{theorem}{Theorem}
\newtheorem{lemma}[theorem]{Lemma}
\newtheorem{proposition}[theorem]{Proposition}
\newcommand{\Z}{\mathbb Z}
\newcommand{\Q}{\mathbb Q}
\newcommand{\C}{\mathbb C}
\newcommand{\HH}{\mathfrak H}
\newcommand{\G}{\Gamma_\theta}
\newcommand{\blog}{https://terrytao.wordpress.com/2026/10/20/an-experimental-ai-math-challenge-using-tenth-discriminant-bounds-to-explore-arguments-ruling-out-siegel-zeroes/}
\title{Could the class-number-one theorem\\have been proved in 1891?\\
\large A reconstruction from Weber's modular theory}
\author{Research note prepared with Codex}
\date{historical-v0.1.0, revised 22 October 2026}
\begin{document}
\maketitle

\begin{abstract}
We investigate a possible route to the imaginary quadratic
class-number-one theorem using the modular theory in Weber's
\emph{Elliptische Functionen und algebraische Zahlen} (1891).
The central deduction is Heegner's: comparing cubic equations for a
singular value and its fourth power produces a Diophantine equation.
We give that deduction, a complete elementary solution of the equation,
and the bound \(q<196\), followed by a four-case exclusion above 163.
We also give a specialization of Stark's repair of Weber's
class-invariant argument, with explicit control of the roots of the
replacement modular equation. The product construction used in this
repair is already part of Weber's 1891 method. This supports 1891 as a
historical upper-bound proposal under a stated rule permitting new
applications of published methods. It neither dates Heegner's discovery
to 1891 nor establishes the earliest possible year. Classical modular
results imported from Weber are stated explicitly and located in the
original edition; the paper does not rebuild the whole theory of
elliptic functions from elementary analysis.
\end{abstract}

\section{The historical question and the scope of the reconstruction}

The question communicated by Ben Green through Terence Tao
\cite{tao} asks how early one could have proved that there is no tenth
imaginary quadratic discriminant of class number one, using results and
concepts already in the literature. It invites a form of historical
research concerned with what may have been possible to prove, alongside
the usual questions of who proved it and when. A reconstruction should
identify both the available methods and the new insight still needed to
connect them.

Pratt \cite{pratt} proposed \(N\le1891\), drawing attention to
Weber's original book \cite{weber}, rather than just the later third
volume of his \emph{Lehrbuch der Algebra} cited by Stark
\cite{stark}. The relevant class invariants and examples occur in
\S\S92--95 of the 1891 edition. Our purpose is to make that proposal
concrete and reviewable, particularly at the step where Weber's proof
requires repair.

We use the following historical criterion. Substantial general results
and organizing methods must have been published by the proposed date.
New applications, deductions for the particular problem, and completions
of an existing argument using its existing method are allowed. Merely
deriving a substantial later method from older foundations does not
backdate that method. This last distinction disqualifies the determinant
route in an earlier reconstruction of this project as an answer to the
present historical question; that route is not used here.

Heegner's passage from modular equations to a Diophantine restriction
\cite{heegner} remains the decisive new observation in the reconstruction.
It is not represented as an observation published in 1891. The case for
historical eligibility is that Weber already developed the particular
class invariants, their degree reductions, and the factorization
calculations to which this observation applies. We return to the limits
of this claim in Section~\ref{assessment}.

\subsection{Mathematical statement}

Write \(D=-q<0\) for a fundamental discriminant, and let \(h(D)\)
be the number of proper equivalence classes of primitive positive
definite forms
\[
 Ax^2+Bxy+Cy^2,\qquad B^2-4AC=D.
\]
Proper equivalence means an integral substitution of determinant one.
The form convention is explained directly below, so no theory of ideals
is required for the arithmetic reductions.

\begin{theorem}\label{main}
Using the classical modular facts stated in Section~\ref{toolkit},
the argument below proves
\[
 h(-q)=1\quad\Longleftrightarrow\quad
 q\in\{3,4,7,8,11,19,43,67,163\}.
\]
For a hypothetical additional discriminant with \(q>163\), it first
gives \(q<196\). All the imported modular facts are in Weber's 1891
book; the class-invariant step requiring repair is proved in
Sections~\ref{repair}--\ref{square-invariant}.
\end{theorem}

The proof is self-contained after the explicitly stated classical
modular toolkit: it supplies the special modular-equation repair, the
algebraic deduction, the Diophantine descent, and the finite arithmetic.
It makes no use of a modern reciprocity theorem, class-field existence,
the general irreducibility of class polynomials, or a theorem on integral
points on elliptic curves. References marked H are historical inputs
published by 1891. References marked M supply later provenance,
interpretation, or expository guidance. These labels do not themselves
establish historical eligibility; the source comparison in
Section~\ref{assessment} does that work.

\section{The elementary form obstruction}\label{forms}

The small-prime obstruction emphasized by Silberman in the current
discussion \cite{silberman} has a direct formulation in the quadratic
forms of Gauss \cite{gauss}. In modern language, a prime that splits or
ramifies must be the norm of an element when the class number is one;
the imaginary part of that element forces the prime to be large. Here
is the form-theoretic argument, including the exceptional even cases.

If \(D=-4m\) is fundamental, then \(m\) is squarefree and
\(m\equiv1,2\pmod4\). Its principal form is \(x^2+my^2\).
For \(m\equiv2\pmod4\), the primitive form \((2,0,m/2)\)
represents 2; for \(m\equiv1\pmod4\), use
\((2,2,(m+1)/2)\). If there is only one class, the principal form
represents 2, which is impossible for \(m>2\). Thus an additional
discriminant has \(q\) odd, squarefree, and \(q\equiv3\pmod4\).

Its principal form is
\begin{equation}\label{principal}
 Q_q(u,v)=u^2+uv+\frac{q+1}{4}v^2
         =(u+v/2)^2+qv^2/4.
\end{equation}
Suppose an odd prime \(p\) satisfies \(B^2\equiv-q\pmod p\).
Choose \(B\) odd. Then
\[
 \left(p,B,\frac{B^2+q}{4p}\right)
\]
is an integral positive definite form representing \(p\). It is
primitive: either \(p\nmid B\), or squarefreeness of \(q\)
implies \(p\nmid(B^2+q)/(4p)\). Class number one forces
\(p=Q_q(u,v)\); since a prime is not a square, \(v\ne0\), and
\begin{equation}\label{obstruction}
 p\ge q/4.
\end{equation}
If \(q\equiv7\pmod8\), the primitive form
\((2,1,(q+1)/8)\) represents 2, giving the same contradiction for
\(q>8\). Therefore
\begin{equation}\label{prime-shape}
 q>163,\ h(-q)=1\quad\Longrightarrow\quad
 q\text{ prime},\quad q\equiv3\pmod8,\quad 3\nmid q.
\end{equation}
Indeed, a composite squarefree \(q>16\) has a prime divisor
\(p\le\sqrt q<q/4\), contrary to \eqref{obstruction}.

\section{The classical modular toolkit}\label{toolkit}

For \(z\in\HH=\{z:\operatorname{Im}z>0\}\), set
\begin{align*}
 f(z)&=e^{-\pi iz/24}\prod_{n\ge1}(1+e^{(2n-1)\pi iz}),\\
 f_1(z)&=e^{-\pi iz/24}\prod_{n\ge1}(1-e^{(2n-1)\pi iz}),\\
 f_2(z)&=\sqrt2\,e^{\pi iz/12}\prod_{n\ge1}(1+e^{2n\pi iz}).
\end{align*}
We use the classical normalization \(j(z)=e^{-2\pi iz}+744+\cdots\),
and the holomorphic cube root
\[
 \gamma_2(z)=\frac{f(z)^{24}-16}{f(z)^8},\qquad \gamma_2(z)^3=j(z).
\]
This branch is real and positive on the positive imaginary axis.
The defining products converge, and \(f\) has no zero in \(\HH\).
Weber's product and transformation identities give
\begin{equation}\label{weber-cubic}
 f(z)^8,\quad -f_1(z)^8,\quad -f_2(z)^8
 \quad\text{as the roots of}\quad X^3-\gamma_2(z)X-16.
\end{equation}

Here are the precise background facts imported from Weber. They are
statements of nineteenth-century modular theory, not assumptions about
the class-number-one classification.

\begin{enumerate}
\item \textbf{Modular functions and transformation laws.}
The function \(j\) separates orbits of
\(\mathrm{SL}_2(\Z)\). The function \(U=f^{24}\) separates orbits of
\[
 \G=\langle S,T^2\rangle,\qquad S(z)=-1/z,\quad T(z)=z+1.
\]
Modulo 2, the matrices of \(\G\) are \(I\) and
\(\left(\begin{smallmatrix}0&1\\1&0\end{smallmatrix}\right)\).
The modular functions for \(\G\) are rational functions of \(U\).
Its two cusps have \(U=0\) and \(U=\infty\), represented by 1 and
\(i\infty\). Consequently a modular function for \(\G\) that is
holomorphic on \(\HH\) and finite at the first cusp is a polynomial
in \(U\), provided it is meromorphic at the second cusp.
These facts are in \S\S49 and 76, pp.~147--150 and 267--268.

For later use put \(F=f^3\), and \(\zeta_8=e^{\pi i/4}\). Then
\begin{equation}\label{basic-transform}
 F(Sz)=F(z),\qquad F(z+2)=\zeta_8^{-1}F(z).
\end{equation}
For \(\gamma=\left(\begin{smallmatrix}a&b\\c&d\end{smallmatrix}\right)
\in\G\), Weber's formula in \S35, p.~107, after cubing, is
\begin{equation}\label{multiplier}
 F(\gamma z)=e^{-\pi i L(\gamma)/8}F(z),\qquad
 L(\gamma)=db+ca+2d^2-2dc-2da\pmod {16}.
\end{equation}
In particular the multiplier is 1 if \(a,d\) are divisible by 16.

\item \textbf{The class polynomial.}
If \(z\) is the upper-half-plane root of a primitive integral
quadratic equation of negative discriminant \(D\), then \(j(z)\)
is an algebraic integer satisfying a monic polynomial in \(\Z[X]\)
of degree \(h(D)\). In particular \(h(D)=1\) implies
\(j(z)\in\Z\). We need only this degree bound, not irreducibility.
The integrality proof is in \S87 and the class-polynomial proof in
\S88, pp.~331--337. The latter uses an induction and polynomial
greatest common divisors, not a class-field theorem.

\item \textbf{The level-two equation.}
There is a polynomial \(\Phi_2(X,Y)\in\Z[X,Y]\), monic of degree
three in \(X\), with
\begin{equation}\label{phi2}
 \Phi_2(X,j(z))=
 (X-j(2z))(X-j(z/2))(X-j((z+1)/2)).
\end{equation}
See \S72, pp.~247--255, and its use in \S89, p.~338.

\item \textbf{The cube-root class invariant.}
If \(z\) satisfies the primitive equation
\(Az^2+Bz+C=0\), where \(3\nmid A(B^2-4AC)\) and \(3\mid B\),
then
\begin{equation}\label{gamma-invariant}
 \gamma_2(z)\in\Q(j(z)).
\end{equation}
This is the theorem of \S92, pp.~350--352, using the transformation
equations of \S74. Its proof selects one of the three cube roots
of \(j(z)\) by the congruence \(3\mid B\), then takes a polynomial
greatest common divisor. This is a different, already established
class invariant from the square invariant discussed next.
\end{enumerate}

Weber's \S94, Theorems I and II, asserts the additional square
invariant needed below. We do \emph{not} simply import its proof:
Sections~\ref{repair}--\ref{square-invariant} supply the missing
root control. The stronger cube-invariant assertion stated
conjecturally on p.~365 is never used. Appendix~\ref{audit}
records the distinction and the normalization of
\eqref{multiplier}.

\section{A modular equation with exactly the required roots}\label{repair}

Stark \cite[\S3]{stark} identified the difficulty with Weber's chosen
Schl\"afli equation: an algebraic relation alone does not ensure that
all its roots come from the intended modular transforms. Additional
roots would invalidate the subsequent greatest-common-divisor
argument. His remedy is to construct the polynomial as a product over
those transforms. The following is a specialization to an odd prime
\(p\), sufficient for our purpose. The construction is the one
described in Weber's \S76, pp.~267--269, with the rationality argument
made explicit.

Define
\[
 M_\infty z=pz,\qquad M_\nu z=\frac{z+16\nu}{p}
       \quad(0\le\nu<p).
\]
For \(e=1\) or \(8\), form
\begin{equation}\label{product}
 \prod_{M\in\{M_\infty,M_0,\ldots,M_{p-1}\}}
 \bigl(X-F(Mz)^eF(z)^{23pe}\bigr)
       =\Sigma_{p,e}(X,U(z)).
\end{equation}

\begin{lemma}\label{product-lemma}
The right side of \eqref{product} is a polynomial
\(\Sigma_{p,e}(X,U)\in\Q[X,U]\). Set
\[
 R_{p,e}(X)=\Sigma_{p,e}(X^{23p+1},X^{8/e}).
\]
This is a nonzero polynomial in \(\Q[X]\). For every
\(w\in\HH\),
\begin{equation}\label{root-test}
 R_{p,e}(F(w)^e)=0
 \quad\Longleftrightarrow\quad
 F(Mw)^e=F(w)^e\text{ for at least one }M.
\end{equation}
\end{lemma}

\begin{proof}
We first verify that the factors in \eqref{product} are permuted by
\(T^2\) and \(S\). Under \(z\mapsto z+2\), choose \(\nu'\)
with \(16\nu'\equiv16\nu+2\pmod p\). Then
\[
 M_\nu(z+2)=M_{\nu'}z+\ell,\qquad
 \ell=(16\nu+2-16\nu')/p\equiv2p\pmod {16}.
\]
Here \(\ell\) is even; the last congruence follows from
\(p^2\equiv1\pmod8\). Thus \(F(M_\nu(z+2))=
\zeta_8^{-p}F(M_{\nu'}z)\). The same multiplier occurs for
\(M_\infty\). After taking the \(e\)-th power, it cancels the multiplier
\(\zeta_8^{-23pe}\) of \(F(z)^{23pe}\), since \(24pe\)
is divisible by 8.

Under \(S\), the transforms \(M_\infty\) and \(M_0\) exchange.
For \(0<\nu<p\), choose \(0<\nu'<p\) with
\(256\nu\nu'\equiv-1\pmod p\). Then
\[
 M_\nu S=\gamma M_{\nu'},\qquad
 \gamma=\begin{pmatrix}
 16\nu&(-1-256\nu\nu')/p\\ p&-16\nu'
 \end{pmatrix}\in\G.
\]
Its diagonal entries are divisible by 16, so \eqref{multiplier}
gives \(F(\gamma z)=F(z)\). This proves the claimed permutation.

The coefficients of the product are therefore modular functions for
\(\G\), holomorphic on \(\HH\). They are finite at the cusp 1:
each \(M(1)\) is a rational number with odd numerator and odd
denominator after reduction, so is equivalent to that cusp, where
\(F\) tends to zero. Their product expansions show that they are
meromorphic at infinity. The first toolkit fact now makes each
coefficient a polynomial in \(U\).

To check its numerical coefficients, put \(Q=e^{\pi iz/8}\). Then
\[
 F(z)=Q^{-1}\prod_{n\ge1}(1+Q^{8(2n-1)})^3
       \in\Z((Q)).
\]
In \(F(M_\nu z)\), replace \(Q\) by
\(\zeta_p^\nu Q^{1/p}\), where \(\zeta_p=e^{2\pi i/p}\).
Every automorphism \(\zeta_p\mapsto\zeta_p^a\) permutes the
factors indexed by \(\nu\); the remaining factor has rational
coefficients. Thus the coefficient series of \eqref{product} are
rational. Invariance under \(T^2\) makes them series in integral
powers of \(Q^8\). Since \(U=Q^{-8}(1+\cdots)\) has rational
coefficients and leading coefficient 1, successive subtraction of
its powers shows that the polynomial coefficients belong to \(\Q\).

Substitute \(X=F(w)^{e(23p+1)}\) in \eqref{product}. Since
\(F(w)\ne0\), factoring out \(F(w)^{23pe}\) from each factor
gives exactly \eqref{root-test}. Finally the powers of \(Q\)
at infinity show that no \(F(Mz)^e\) is identically \(F(z)^e\):
their leading exponents are multiplied by \(p\) or \(1/p\).
Their finite product of differences is not identically zero, so
\(R_{p,e}\) is not the zero polynomial.
\end{proof}

Only rational coefficients are needed. The construction in fact permits
an integrality refinement, but it contributes nothing to the proof.
The purpose of the product is the equivalence \eqref{root-test},
which rules out precisely the extra roots at issue in Weber's proof.

\section{Recovering the square class invariant}\label{square-invariant}

From now on let \(q\) satisfy \eqref{prime-shape}, and put
\[
 \theta=i\sqrt q,\qquad J=j(\theta).
\]
We prove the special case of Weber's assertion required here:
\begin{equation}\label{f2-invariant}
 f(\theta)^2\in\Q(J).
\end{equation}

\subsection{The twenty-fourth power}

By \eqref{weber-cubic}, the three roots of
\begin{equation}\label{U-cubic}
 (X-16)^3-JX=0
\end{equation}
are
\[
 u_0=f(\theta)^{24},\quad
 u_1=f(\theta+1)^{24},\quad
 u_2=f(1-1/\theta)^{24}.
\]
We claim that only \(u_0\) is a root of \(R_{q,8}\).
It is a root because \(M_0\theta=\theta/q=-1/\theta\).

For the exclusion, suppose \(w\) satisfies a primitive equation
\(Aw^2+2Bw+C=0\) with \(AC-B^2=q\), and suppose
\(f(Mw)^{24}=f(w)^{24}\). The orbit-separation property gives
\(Mw=\gamma w\) for \(\gamma\in\G\). The integral matrix
\(\gamma^{-1}M\), of determinant \(q\), fixes \(w\), so is
of the form
\begin{equation}\label{fixer}
 \begin{pmatrix}h-B\ell&-C\ell\\A\ell&h+B\ell\end{pmatrix},
 \qquad h,\ell\in\Z,\qquad h^2+q\ell^2=q.
\end{equation}
The integrality of \(\ell\) follows by comparing the fixed-point
quadratic with the primitive equation; then \(h\) is integral too.
As \(q\) is not a square, necessarily \(h=0,\ell=\pm1\).
Every \(M\) is the identity modulo 2. Consequently the matrix in
\eqref{fixer}, modulo 2, must be either \(I\) or
\(\left(\begin{smallmatrix}0&1\\1&0\end{smallmatrix}\right)\).

For \(w=\theta+1\), its coefficients are
\((A,B,C)=(1,-1,q+1)\); for \(w=1-1/\theta\), they are
\((q,-q,q+1)\). In either case the matrix is
\(\left(\begin{smallmatrix}1&0\\1&1\end{smallmatrix}\right)\)
modulo 2, which is forbidden. The claim follows from
\eqref{root-test}.

The root \(u_0\) is simple in \eqref{U-cubic}. For example, its
derivative there is
\((u_0-16)^2(2+16/u_0)>0\); the product gives \(u_0>16\)
in our range. Taking the monic greatest common divisor of
\eqref{U-cubic} and \(R_{q,8}\) over \(\Q(J)\) therefore gives
\(X-u_0\). Hence
\begin{equation}\label{f24-field}
 f(\theta)^{24}\in\Q(J).
\end{equation}

\subsection{The sixth power}

The eight roots of \(X^8-f(\theta)^{24}\) are
\[
 F(\theta+2j)=\zeta_8^{-j}F(\theta),\qquad 0\le j<8.
\]
We determine which are also roots of \(R_{q,1}\). Write
\(w=\theta+2j\), \(B=-2j\), so that
\(w^2+2Bw+(q+B^2)=0\). As above, equality of the \(F\) values
first implies equality of the \(f^{24}\) values. Formula
\eqref{fixer} then applies with \(h=0,\ell=\pm1\).
The transform \(M_\infty\) is impossible because the lower-left
entry of the resulting \(\gamma\) would be \(\pm1/q\).
For \(M_\nu\), integrality forces \(16\nu\equiv B\pmod q\).
For this unique \(\nu\), up to an irrelevant overall sign,
\[
 \gamma=\begin{pmatrix}a&1+Ba\\-1&-B\end{pmatrix},
 \qquad a=(B-16\nu)/q.
\]
Since \(B\) is even and \(q\equiv3\pmod8\),
\(a\equiv3B\pmod {16}\). Substitution in
\eqref{multiplier} gives
\[
 L(\gamma)=-3B+2B^2-a(B-1)^2
       \equiv6B\equiv-12j\pmod {16}.
\]
For clarity, the last congruence follows on putting \(B=2b\):
\(L-6B\) is congruent to
\(-8b(3-4b+3b^2)\), divisible by 16.
Thus the multiplier is \((-i)^j\), and equality of the \(F\)
values holds exactly for \(j=0,4\). These give
\(F(\theta)\) and \(-F(\theta)\).

The monic greatest common divisor of \(X^8-f(\theta)^{24}\)
and \(R_{q,1}\) is consequently
\(X^2-f(\theta)^6\). Together with \eqref{f24-field}, this proves
\begin{equation}\label{f6-field}
 f(\theta)^6\in\Q(J).
\end{equation}

Finally, the cube-root class-invariant theorem
\eqref{gamma-invariant} applies to \(\theta^2+q=0\).
It gives \(\gamma_2(\theta)\in\Q(J)\), and this value is nonzero
by its product expression in our range. Therefore
\[
 f(\theta)^8=\frac{f(\theta)^{24}-16}{\gamma_2(\theta)}\in\Q(J),
 \qquad
 f(\theta)^2=\frac{f(\theta)^8}{f(\theta)^6}\in\Q(J).
\]
This proves \eqref{f2-invariant}. It is the combination of Weber's
eighth- and sixth-power results on pp.~360--362; it uses no conjectural
statement about \(f(\theta)\) itself.

\section{Heegner's cubic comparison}\label{comparison}

The step in this section is Heegner's \cite{heegner}, as explained by
Stark \cite[\S2]{stark}. We include every algebraic calculation and use
a rational-root argument to avoid the general class-polynomial
irreducibility theorem.

Put
\[
 \tau=\frac{3+\theta}{2},\qquad
 g=\gamma_2(\tau),\qquad \alpha=\frac{2}{f(\theta)^2}.
\]
The equation of \(\tau\) is
\(\tau^2-3\tau+(q+9)/4=0\), of discriminant \(-q\).
The class polynomial gives \(j(\tau)\in\Z\), and
\eqref{gamma-invariant} gives \(g\in\Q\). As
\(g^3=j(\tau)\), the rational-root theorem gives \(g\in\Z\).
Since \(2\tau=3+\theta\), \eqref{phi2} gives
\([\Q(J):\Q]\le3\). In view of \eqref{f2-invariant},
\begin{equation}\label{degree-upper}
 [\Q(\alpha):\Q]\le3.
\end{equation}

Weber's product identity in \S29, used explicitly in \S95,
pp.~367--368, gives
\begin{equation}\label{alpha-normalization}
 \alpha=e^{-\pi i/4} f_2(\tau)^2.
\end{equation}
It can also be checked directly: for \(t=e^{\pi i\theta}\), separating
even and odd indices gives
\[
 \prod_{n\ge1}(1+(-t)^n)\prod_{n\ge1}(1+t^{2n-1})=1.
\]
Substitution in the definitions of \(f\) and \(f_2\) gives exactly
\eqref{alpha-normalization}, including its phase and factor 2.
In particular \(\alpha^4=-f_2(\tau)^8\); applying
\eqref{weber-cubic} at \(\tau\) yields
\begin{equation}\label{degree12}
 \alpha^{12}-g\alpha^4-16=0.
\end{equation}
Thus \(\alpha\) is an algebraic integer. The factor 2 in its
definition and the phase in \eqref{alpha-normalization} are both
essential. The product at the purely imaginary point \(\theta\)
also gives
\begin{equation}\label{alpha-small}
 0<\alpha<2e^{-\pi\sqrt q/12}<1.
\end{equation}

Set \(t=\alpha^4\). Then \(0<t<1\) and
\[
 P(X)=X^3-gX-16,\qquad g=t^2-16/t<0.
\]
The polynomial \(P\) is strictly increasing on the real line, so its
only real root is \(t\in(0,1)\). A reducible monic integral cubic
has an integer root; \(P\) has none. Hence it is irreducible, and
\begin{equation}\label{degree-exact}
 [\Q(\alpha^4):\Q]=[\Q(\alpha):\Q]=3.
\end{equation}
This proves the particular degree statement needed without importing
a theorem that all singular moduli have degree equal to their class
number.

Choose \(\beta=\alpha\) or \(-\alpha\), and write its monic
minimal polynomial as
\[
 X^3+aX^2+bX+c,\qquad a,b,c\in\Z.
\]
Squaring the roots produces a cubic with coefficients
\[
 e=2b-a^2,\qquad f=b^2-2ac,\qquad k=-c^2.
\]
Squaring again produces
\[
 X^3+(2f-e^2)X^2+(f^2-2ek)X-k^2.
\]
By \eqref{degree-exact}, this is the minimal polynomial \(P\) of
\(\beta^4=\alpha^4\). Comparing constant terms gives \(c^4=16\).
Change the sign of \(\beta\) if necessary so that \(c=2\).
The coefficient of \(X^2\) then gives
\begin{equation}\label{coefficient-relation}
 2(b^2-4a)=(2b-a^2)^2.
\end{equation}
Modulo 2, this implies \(a\) even; modulo 4 it then implies
\(b\) even. Set
\[
 x=-a/2,\qquad y=(b-a^2)/2.
\]
Equation \eqref{coefficient-relation} becomes
\begin{equation}\label{heegner-equation}
 y^2=2x(x^3+1).
\end{equation}
Pratt's version \(y^2=2x(x^3-1)\) is the same equation after
\(x\mapsto-x\).

\section{Solving the Diophantine equation}\label{diophantine}

\begin{lemma}\label{solutions}
The integral solutions of \eqref{heegner-equation} are
\[
 (x,y)=(0,0),\quad(-1,0),\quad(1,\pm2),\quad(2,\pm6).
\]
\end{lemma}

\begin{proof}
Deal first with \(x=0,-1\). Otherwise
\(\gcd(x,x^3+1)=1\); in the product \(2x(x^3+1)\), each
odd prime has an even exponent. The signed squarefree part of
\(x^3+1\) is therefore one of \(1,-1,2,-2\), giving four cases.

If \(x^3+1=z^2\), Euler's descent gives \(x=-1,0,2\).
Euler published this result in 1747 \cite[Theorem 10]{euler};
Appendix~\ref{descent} supplies the descent, including the decrease
and the terminal case.

If \(x^3+1=-z^2\), put \(m=-x\), so \(m^3=z^2+1\).
If \(z\) is odd, this would be \(m^3\equiv2\pmod8\),
impossible. Otherwise \(m\) is odd. In \(\Z[i]\) the factors
\(z+i,z-i\) are coprime: a common prime would divide \(2i\)
and the odd norm \(m^3\). Unique factorization therefore gives
\(z+i=(a+bi)^3\); units can be absorbed into cubes. Imaginary
parts give \(b(3a^2-b^2)=1\), forcing \(a=0,b=-1\).
Thus \(z=0,m=1\), the already treated solution \(x=-1\).

If \(x^3+1=2z^2\), equation \eqref{heegner-equation} forces
\(x=w^2\) for an integer \(w\). Then
\[
 (w^2+1)(w^4-w^2+1)=2z^2.
\]
The factors are coprime: a common divisor divides 3, but 3 never
divides \(w^2+1\). The second factor is odd, so must be a square.
For \(|w|\ge2\), however,
\[
 (w^2-1)^2<w^4-w^2+1<(w^2)^2.
\]
The value \(w=0\) fails, and \(w=\pm1\) gives \(x=1\).

If \(x^3+1=-2z^2\), put \(m=-x\). Now
\(m^3=1+2z^2\) is odd. The factors
\(1\pm z\sqrt{-2}\) are coprime in \(\Z[\sqrt{-2}]\):
a common prime would divide 2 and an odd norm. Unique factorization
gives \(1+z\sqrt{-2}=(a+b\sqrt{-2})^3\), since either unit
is a cube. Real parts yield
\(a(a^2-6b^2)=1\), forcing \(a=1,b=0\). Again \(x=-1\).

For completeness, the two unique-factorization assertions here need
only the Euclidean algorithm in these specific rings. Given a quotient
in \(\Q(i)\) or \(\Q(\sqrt{-2})\), round its two coefficients to
integers. The error has norm at most \(1/2\) or \(3/4\),
respectively. Thus division has a remainder of strictly smaller norm.
The Euclidean algorithm supplies greatest common divisors and Bezout
identities; hence irreducibles are prime. Factoring by decreasing
positive integer norms proves existence and uniqueness. The units are
\(\{\pm1,\pm i\}\) and \(\{\pm1\}\). These are explicit
Euclidean factorizations, with no invocation of ideal class theory.
Substituting the resulting values of \(x\) gives the asserted list.
\end{proof}

Reversing the substitution in Section~\ref{comparison}, the possible
pairs of cubic coefficients are
\begin{equation}\label{coefficient-list}
 (a,b)=(0,0),(2,4),(-2,0),(-2,8),(-4,4),(-4,28).
\end{equation}
In particular, \(|a|\le4\) and \(0\le b\le28\).

\section{The bound and the finite finish}\label{finish}

The cubic equation of \(\beta=\pm\alpha\), with constant term 2,
and \eqref{coefficient-list} imply
\[
 2=|\beta^3+a\beta^2+b\beta|
 \le\alpha^3+4\alpha^2+28\alpha<33\alpha.
\]
Together with \eqref{alpha-small} this gives
\[
 e^{\pi\sqrt q/12}<33,
 \qquad \sqrt q<\frac{12\log33}{\pi}<14,
 \qquad q<196.
\]
For the last inequality it is enough to use \(\pi>3\) and
\(e^{7/2}>33\). If an exact arithmetic check of the latter is desired,
the exponential series through degree 11 is
\[
 \sum_{n=0}^{11}\frac{(7/2)^n}{n!}
       =\frac{386627544403}{11678515200}>33.
\]

Only \(171,179,187,195\) lie strictly between 163 and 196 and are
3 modulo 8. Each has a witness to the obstruction
\eqref{obstruction}:
\[
\begin{array}{c|rrrr}
 q&171&179&187&195\\ \hline
 p&3&3&11&3\\
 r&0&1&0&0
\end{array}
\qquad p\mid q+r^2,\quad 4p<q.
\]
This contradicts class number one in every case. One may instead
discard the three composite values immediately and use only the
witness \(p=3\) for 179. Appendix~\ref{finite} gives the small
remaining classification below 164 and checks the nine survivors by
reduction. No singular modulus has to be evaluated numerically.

\section{What the date 1891 does and does not mean}\label{assessment}

The evidence for this cutoff is more specific than a list of old
foundations. Weber's \S76, pp.~267--269, explicitly constructs
transformation equations by taking functions permuted under
\(z\mapsto z+2\) and \(z\mapsto-1/z\), and expressing their
symmetric functions through \(f^{24}\). His \S94 gives the
eighth- and sixth-power invariants and explicitly combines them to
obtain the square invariant. Section~\ref{repair} repairs the root
control by staying within that same product construction. The
cyclotomic coefficient argument uses the familiar action of
automorphisms on roots of unity; it is not a general reciprocity law.

There is also direct evidence for the final algebraic setting.
On p.~368 Weber factors the equations for the known examples through
degrees 24, 12, 6, and 3, and compares their coefficients. What he does
not do is turn the coefficient restrictions into a Diophantine equation
that exhausts all possibilities. That is Heegner's observation. Stark's
retrospective assessment \cite[\S5]{stark} is that Weber could have
completed the argument had he noticed this consequence. Pratt's
bibliographical contribution \cite{pratt} is to locate the material
in 1891, rather than stopping at the later edition Stark cited.

Under the criterion in Section~1, we therefore regard
\(N\le1891\) as a defensible proposal: the substantial modular
methods are already present, the repair is an application of one of
them, and the remaining new step concerns the particular
class-number-one problem. This is a judgment about admissibility,
supported by the displayed construction. If the criterion were to
require every application and every repair itself to have appeared in
print by the cutoff, this reconstruction would not satisfy it. That
stricter rule would answer a different question from whether a proof
could have been assembled using the available methods.

Nor does the argument show that discovery in 1891 was likely. A possible
reason for the delay is that calculating striking singular values and
proving a uniform classification suggest different directions of work:
Weber starts with known discriminants and derives their small equations;
Heegner uses the necessary form of those equations to rule out unknown
discriminants. This reversal is visible in the mathematics, but an
explanation of why a particular mathematician did not make it would
require evidence from that mathematician's work and correspondence.
We offer the contrast as a question for historical research, not as an
estimated probability of discovery.

Finally, no earliest year is proved. The present source audit identifies
1891 as a documented point at which the needed material is together.
It does not rule out an earlier route or an earlier publication of
enough of Weber's ingredients. The credit for the actual theorem and
its later clarification remains with Heegner and the mathematicians
who completed and established its proofs.

\appendix
\section{Euler's descent, with the decrease made explicit}\label{descent}

Euler's Theorem 10 in \cite{euler}, written in 1738 and published
in 1747, treats a cube plus one that is a square. We reproduce the
descent in a form that also covers zero and negative rational inputs.
Conrad's exposition \cite{conrad} is useful guidance for making the
sign choices, the descent measure, and the terminal case explicit.
These details elaborate Euler's factorization and rational
parametrizations; they use no theory of elliptic curves.

\begin{lemma}\label{euler}
The rational solutions of \(Y^2=X^3+1\) are
\[
 (-1,0),\quad(0,\pm1),\quad(2,\pm3).
\]
\end{lemma}

\begin{proof}
First establish the descent assertion: if \(u,v\) are coprime
positive integers, \(3\nmid u\), and
\begin{equation}\label{descent-condition}
 u,\ v,\ u^2-3uv+3v^2\quad\text{are squares},
\end{equation}
then \(u=v=1\).

Write \(w^2=u^2-3uv+3v^2\), choosing the sign so that
\(w\equiv-u\pmod3\). Put
\[
 \frac{w-u}{v}=\frac mn,\qquad \gcd(m,n)=1,\quad n>0.
\]
Here \(3\nmid m\) and \(n\mid v\). Squaring
\(w=u+(m/n)v\) and simplifying gives
\[
 \frac uv=\frac{3n^2-m^2}{n(2m+3n)}.
\]
The fraction on the right is reduced. A prime common to
\(3n^2-m^2\) and \(2m+3n\) would divide both \(m\) and
\(3n\), contradicting \(\gcd(m,n)=1\) and \(3\nmid m\);
the other denominator factor \(n\) is also coprime to the numerator.
Since \(u\) is a square not divisible by 3, comparing signs modulo 3
forces
\begin{equation}\label{mn-pair}
 u=m^2-3n^2,\qquad v=-n(2m+3n).
\end{equation}
Choose \(k\) with \(k^2=u\) and \(k\equiv-m\pmod3\), and
reduce the nonzero fraction
\[
 \frac{k-m}{n}=\frac{u'}{v'},\qquad \gcd(u',v')=1,\quad v'>0.
\]
It has \(3\nmid u'\). From
\(m^2-3n^2=(m+(u'/v')n)^2\), one obtains
\begin{equation}\label{new-pair}
 v=n^2\frac{u'^2-3u'v'+3v'^2}{u'v'}.
\end{equation}
The quadratic in the numerator is positive. Since \(v>0\) and
\(v'>0\), also \(u'>0\). The three positive integers
\(u',v',u'^2-3u'v'+3v'^2\) are pairwise coprime: the only
possible additional common prime would be 3 dividing \(u'\).
As \(v\) is a square, multiplying \eqref{new-pair} by
\((u'v'/n)^2\) shows that their product is a rational square,
hence an integer square. Each is therefore a square. This constructs
another pair satisfying \eqref{descent-condition}.

Moreover, \(v/n\) is an integer by \eqref{mn-pair}. Equation
\eqref{new-pair}, together with the coprimality just proved, implies
\(u'v'\mid n\). Thus
\[
 0<v'\le u'v'\le n\le v.
\]
If \(v'=v\), every inequality is equality, so
\(u'=1\), \(v'=n=v\). Then \((k-m)/n=1/n\), or
\(k=m+1\). Combining \(m^2-3n^2=k^2\) with
\(v=n=-n(2m+3n)\) gives
\(2m+1=-3n^2=-3n\), hence \(n=1,m=-2,u=v=1\).
If instead the new pair is already \((u',v')=(1,1)\), then
\(k=m+n\), so \(m=-2n\). Coprimality of \(m,n\) again
forces \(n=1\) and \(u=v=1\).
Consequently any nontrivial pair descends to a nontrivial pair with
strictly smaller positive \(v\), which is impossible. This proves
the descent assertion.

Now let \(Y^2=X^3+1\). Necessarily \(X\ge-1\), and
\(X=-1\) gives \(Y=0\). Otherwise write
\(X=a/b\), \(b>0\), \(\gcd(a,b)=1\), and set
\(c=a+b>0\). The equation implies that
\[
 b(a^3+b^3)=bc(c^2-3bc+3b^2)
\]
is an integer square. If \(3\nmid c\), the three factors are
pairwise coprime and positive, so are squares. Apply the descent
assertion with \((u,v)=(c,b)\) to get \(b=c=1\), and thus
\(X=0\).

If \(c=3d\), then \(3\nmid b\), and after removing a factor 9
the square condition is
\[
 bd(b^2-3bd+3d^2)\text{ is a square}.
\]
The factors are again pairwise coprime. The descent assertion with
\((u,v)=(b,d)\) gives \(b=d=1\), hence \(c=3,a=2,X=2\).
The corresponding values of \(Y\) complete the list.
\end{proof}

\section{The small discriminants and the nine survivors}\label{finite}

We supply the classification below 164 rather than referring to an
earlier paper. Section~\ref{forms} excludes every even fundamental
discriminant except \(-4,-8\), and for odd \(q>8\) forces
\(q\equiv3\pmod8\). The value 15 is excluded by the ramified
prime 3; every composite \(q>16\) is excluded as in
\eqref{prime-shape}. The primes between 9 and 163 congruent to 3
modulo 8 are
\[
 11,19,43,59,67,83,107,131,139,163.
\]
The four values \(59,83,107,131\) have \(-q\equiv1\pmod3\),
and \(139\) has \(-139\equiv1\pmod5\). The primes 3 and 5
are below \(q/4\), excluding these values by
\eqref{obstruction}. Including the smaller fundamental discriminants
leaves exactly the nine in Theorem~\ref{main}.

To check that they do have class number one, use elementary reduction
of positive definite forms \cite{gauss}. Replacing \(B\) by its
residue modulo \(2A\) makes \(|B|\le A\); if \(C<A\), the
substitution \((x,y)\mapsto(-y,x)\) exchanges the outer
coefficients and strictly decreases the positive first coefficient.
Iteration therefore gives
\[
 |B|\le A\le C,\qquad 3A^2\le4AC-B^2=q.
\]
On the boundary, take \(B\ge0\) when \(|B|=A\) or \(A=C\).
The usual substitutions identify the two signs there. For each of the
nine values, the finite check of
\(B^2\equiv-q\pmod{4A}\) gives the following complete table:
\begin{center}
\begin{tabular}{rrl}
\toprule
\(q\)&\(\lfloor\sqrt{q/3}\rfloor\)&Reduced forms \((A,B,C)\)\\
\midrule
3&1&\((1,1,1)\)\\
4&1&\((1,0,1)\)\\
7&1&\((1,1,2)\)\\
8&1&\((1,0,2)\)\\
11&1&\((1,1,3)\)\\
19&2&\((1,1,5)\)\\
43&3&\((1,1,11)\)\\
67&4&\((1,1,17)\)\\
163&7&\((1,1,41)\)\\
\bottomrule
\end{tabular}
\end{center}
Every form is equivalent to a reduced form, and each row has just one.
This verifies the other implication of Theorem~\ref{main}. The finite
work consists only of these small divisions and congruences.

\clearpage
\section{Source audit and the remaining historical judgment}\label{audit}

All page references below are printed pages of the 1891 edition, not
PDF page numbers or page numbers in the 1908 edition.

\begin{center}
\small
\begin{tabular}{>{\raggedright\arraybackslash}p{0.30\textwidth}
 >{\raggedright\arraybackslash}p{0.23\textwidth}
 >{\raggedright\arraybackslash}p{0.38\textwidth}}
\toprule
Input or method&Weber (1891)&Role in this paper\\
\midrule
Product definitions and identities&\S29; \S49,
pp.~149--150&Defines the functions and the cubic
\eqref{weber-cubic}.\\[4pt]
Linear transformation of \(f\)&\S35, pp.~105--107&After cubing,
gives the exact multiplier \eqref{multiplier}.\\[4pt]
Functions for the theta subgroup&\S49, pp.~147--149;
\S76, pp.~267--268&Converts invariant coefficient functions into
polynomials in \(f^{24}\).\\[4pt]
Invariant modular equations&\S72, pp.~247--255&Supplies the
degree-three level-two equation.\\[4pt]
Products over modular transforms&\S\S73,76, pp.~255--269&The
existing organizing method used in the root-control repair.\\[4pt]
Integrality and the class polynomial&\S\S87--88,
pp.~331--337&Only degree at most the class number is needed.\\[4pt]
\(\gamma_2\) as a class invariant&\S92,
pp.~350--352&Applied separately at \(\theta\) and \(\tau\).\\[4pt]
\(f^8,f^6\), and their quotient \(f^2\)&\S94,
pp.~360--362&Statements located; the special case needed is proved
in Section~\ref{square-invariant} with the repaired equation.\\[4pt]
Stronger \(f^3\) assertion&\S94, p.~365&Explicitly conjectural
there; not used.\\[4pt]
Low-degree factorizations for known examples&\S95, pp.~367--368&
Evidence that the specific algebraic setting was already under study.\\
\bottomrule
\end{tabular}
\end{center}

\paragraph{The multiplier convention.}
Weber writes the transform as
\((\gamma+\delta z)/(\alpha+\beta z)\). His final formula
on p.~107 has a cube-root-of-unity factor, a factor
\((2/(\alpha-\beta))\), and the exponential
\(\exp[-3\pi i(\alpha-\beta)(\alpha+\beta+\gamma-\delta)/8]\).
Upon cubing, the first factor disappears. Using
\((2/r)=(-1)^{(r^2-1)/8}\) for odd \(r\) and
\(\alpha\delta-\beta\gamma=1\) gives
\eqref{multiplier} after
\((a,b,c,d)=(\delta,\gamma,\beta,\alpha)\).
This also checks \(S\) and \(T^2\) directly and avoids confusing
Weber's matrix convention with the modern one.

\paragraph{What was actually checked.}
The 1891 statements and the relevant transformation and product-method
passages were compared with the original scan. Stark's full 1969 paper,
especially pp.~20--25, was used to identify the root-control gap and
its repair. The specialization above gives explicit transformations
and congruences for an odd prime, rather than treating Stark's later
paper as a historical black box. Euler's Theorem 10 was located in
the original work and its German translation; its publication date
1747 is recorded in the Euler Archive. The descent above supplies
the details needed for a complete proof rather than relying on
Euler's brief assertion that the new numbers are smaller.

\paragraph{Verification status.}
The accompanying script checks polynomial identities, multiplier
congruences, descent identities, the rational exponential inequality,
and the finite form enumeration using exact arithmetic. These checks
do not verify the analytic toolkit or the historical interpretation.
This working draft has not been assigned an independent audit verdict
and is not a formal verification. The remaining historical judgment
is whether the explicitly displayed repair and Heegner's
problem-specific deduction meet the criterion in Section~1;
the date 1891 is advanced with that criterion visible.

\clearpage
\begin{thebibliography}{M6}

\bibitem[H1]{gauss}
C.~F.~Gauss, \emph{Disquisitiones arithmeticae}, Leipzig, 1801,
Section V, especially Articles 153--159 and 171--175.
\href{https://doi.org/10.3931/e-rara-4025}{Original edition}.

\bibitem[H2]{euler}
L.~Euler, \emph{Theorematum quorundam arithmeticorum
demonstrationes}, Commentarii academiae scientiarum Petropolitanae
\textbf{10} (published 1747; written 1738), 125--146,
Theorem 10, Enestr\"om E98.
\href{https://scholarlycommons.pacific.edu/euler-works/98/}{Euler Archive record and translations};
\href{https://math.dartmouth.edu/~euler/docs/originals/E098.pdf}{original scan}.

\bibitem[H3]{weber}
H.~Weber, \emph{Elliptische Functionen und algebraische Zahlen},
Friedrich Vieweg und Sohn, Braunschweig, 1891.
\href{https://archive.org/details/ellipfuncundalg00weberich}{Original edition}.
This is the edition used for all section and page references here.

\bibitem[M1]{heegner}
K.~Heegner, \emph{Diophantische Analysis und Modulfunktionen},
Mathematische Zeitschrift \textbf{56} (1952), 227--253.
\href{https://doi.org/10.1007/BF01174749}{doi:10.1007/BF01174749}.

\bibitem[M2]{stark}
H.~M.~Stark, \emph{On the ``gap'' in a theorem of Heegner},
Journal of Number Theory \textbf{1} (1969), 16--27.
\href{https://doi.org/10.1016/0022-314X(69)90023-7}{doi:10.1016/0022-314X(69)90023-7};
\href{https://hdl.handle.net/2027.42/33039}{University of Michigan repository}.

\bibitem[M3]{pratt}
K.~Pratt, Comment proposing the historical cutoff \(N\le1891\),
22 October 2026, 8:46 am,
\href{\blog comment-page-1/\#comment-698625}{comment 698625 in Tao's discussion}.

\bibitem[M4]{tao}
T.~Tao, Historical variant communicated by Ben Green,
21 October 2026, 1:10 pm,
\href{\blog\#comment-698559}{comment 698559}
on \emph{An experimental AI-math challenge: using tenth discriminant
bounds to explore arguments ruling out Siegel zeroes}.

\bibitem[M5]{conrad}
K.~Conrad, \emph{An example of descent by Euler}, expository note,
\href{https://kconrad.math.uconn.edu/blurbs/ugradnumthy/descentbyeuler.pdf}{author's PDF}.

\bibitem[M6]{silberman}
L.~Silberman, Comment on small-prime inertness,
20 October 2026, 8:49 am,
\href{\blog\#comment-698449}{comment 698449 in Tao's discussion}.

\end{thebibliography}
\end{document}
